\documentclass[11pt]{article}
\usepackage[utf8]{inputenc}
\usepackage[T1]{fontenc}
\usepackage{amsmath,amssymb,amsfonts}
\usepackage{booktabs}
\usepackage{multirow}
\usepackage{graphicx}
\usepackage{hyperref}
\usepackage[margin=1in]{geometry}
\usepackage{natbib}
\usepackage{xcolor}
\usepackage{array}

\newcommand{\Kstar}{K^{*}}

\title{Recursive Generation Has a Universal Discrimination Boundary:\\
  $\Kstar\!=\!2$ Across Six Architecture Families at 1B--8B Scale}

\author{
  Anonymous Authors
}

\date{}

\begin{document}
\maketitle

% ============================================================
\begin{abstract}
Web-scale corpora are increasingly contaminated by recursively generated synthetic text, yet current detectors treat synthetic text as a binary category and cannot estimate how many rounds of recursive generation a passage has undergone.
We formalize recursive generation as an iterated lossy channel and derive, via the Data Processing Inequality (DPI), that a finite empirical discrimination boundary $\Kstar$ must exist beyond which successive generations become statistically indistinguishable.
Through a systematic experimental matrix spanning 18 conditions---6 model architecture families (Pythia, OLMo, Qwen, Qwen-7B, LLaMA-3.1-8B, Mistral-7B), 8 decoding configurations across 2 models, 3 text lengths, and a $4\!\times\!4$ cross-scorer matrix---we find that $\Kstar\!=\!2$ holds in 17 of 18 conditions (94\%).
The single exception arises in a degenerate self-scorer configuration and has a clear mechanistic explanation.
Two-way ANOVA confirms that the generator identity accounts for 68\% of variance in detection AUC while the scorer identity is not significant ($p\!=\!0.23$), establishing that depth signal is predominantly text-intrinsic and scorer-independent.
$\Kstar\!=\!2$ correctly predicts that binary filtering at the discrimination boundary outperforms graduated depth weighting, surprisal-based selection, and distribution matching for downstream task performance.
\end{abstract}

% ============================================================
\section{Introduction}
\label{sec:intro}

The proliferation of large language models (LLMs) has led to a growing concern: web-crawled training corpora increasingly contain synthetic text that was itself generated by LLMs, potentially trained on earlier synthetic outputs~\citep{shumailov2024model,alemohammad2024selfconsuming}.
This recursive contamination creates a \emph{generational depth} structure---generation~0 is human-written, generation~1 is model-generated from human prompts, generation~2 is model-generated from generation~1 outputs, and so on.

Existing detection methods treat this as a binary problem: human versus machine~\citep{mitchell2023detectgpt,bao2024fastdetectgpt}.
Data selection methods such as DSIR~\citep{xie2023dsir} perform distribution matching without modeling recursive depth.
Model collapse literature~\citep{shumailov2024model,dohmatob2024tale} characterizes macro-level degradation but does not provide per-document depth estimation tools.
\citet{drayson2025contamination} applies binary machine-generated text detection to importance resampling, the closest prior work, but still makes binary judgments rather than estimating recursive depth.

We ask: \emph{how many generations of recursive synthesis can a statistical detector reliably distinguish?}
We formalize this question through the Data Processing Inequality (DPI).
Treating each generation step as a noisy channel, DPI guarantees that mutual information between the original text and generation~$k$ is non-increasing in $k$.
This implies a finite \emph{discrimination boundary} $\Kstar_{\mathrm{theory}}$: the depth beyond which no detector can reliably separate adjacent generations.
We introduce $\Kstar_{\mathrm{empirical}}$, a conservative lower bound on $\Kstar_{\mathrm{theory}}$ measured using 15 hand-crafted statistical features and a LightGBM classifier with bootstrap confidence intervals.

Our key finding is that $\Kstar\!=\!2$ is \emph{universal}: it holds across 6 architecture families (Pythia, OLMo, Qwen-1.5B, Qwen-7B, LLaMA-3.1-8B, Mistral-7B-v0.3), spanning 1B to 8B parameters, across 8 decoding configurations (greedy, nucleus $p\!\in\!\{0.9, 0.95\}$, temperature $T\!\in\!\{0.7, 1.0, 1.2\}$) tested on 2 model families, across text lengths from 128 to 512 tokens, and across a $4\!\times\!4$ generator$\times$scorer matrix.
In total, 17 of 18 tested conditions yield $\Kstar\!=\!2$, with the single exception attributable to a degenerate self-scorer configuration where the weakest-signal generator scores itself.

\paragraph{Contributions.}
(1)~We formalize $\Kstar$ as an operational quantity derived from DPI and provide the first systematic empirical measurement.
(2)~We demonstrate $\Kstar\!=\!2$ universality across 6 architecture families at 1B--8B scale, 8 decoding configurations, 3 text lengths, and a $4\!\times\!4$ cross-scorer matrix (17/18 conditions).
(3)~We show that $\Kstar\!=\!2$ serves as an actionable data curation guideline: binary filtering at $\Kstar$ outperforms graduated weighting and alternative selection methods.

% ============================================================
\section{Related Work}
\label{sec:related}

\paragraph{Machine-generated text detection.}
Statistical detectors leverage perplexity~\citep{solaiman2019release}, log-rank perturbation~\citep{mitchell2023detectgpt}, and curvature-based metrics~\citep{bao2024fastdetectgpt} for binary classification.
These methods achieve high accuracy for generation-0 versus generation-1 but do not model recursive depth.

\paragraph{Model collapse and iterative training.}
\citet{shumailov2024model} demonstrate that training on self-generated data leads to progressive quality degradation.
\citet{alemohammad2024selfconsuming} and \citet{dohmatob2024tale} provide theoretical frameworks for understanding this collapse.
Our work complements this line by characterizing the \emph{detection-side} boundary rather than the training-side dynamics.

\paragraph{Data selection and curation.}
DSIR~\citep{xie2023dsir}, D4~\citep{tirumala2024d4}, and DoReMi~\citep{xie2024doremi} optimize training mixtures via distribution matching or domain reweighting.
\citet{drayson2025contamination} applies binary MGT detection to importance resampling for preventing model collapse.
None of these methods estimate recursive generation depth or provide a principled boundary for when depth-aware filtering adds value over binary detection.

% ============================================================
\section{Method}
\label{sec:method}

\subsection{Recursive Generation as an Iterated Lossy Channel}
\label{sec:dpi}

Let $X_0$ denote human-written text drawn from distribution $P_0$.
A language model $M$ defines a stochastic channel that maps $X_k$ to $X_{k+1}$: the model reads $X_k$ (or a prompt derived from it) and generates a continuation or rewrite.
We consider continuation mode as our canonical condition, where $M$ is given the first $n$ tokens of $X_k$ and generates a completion.

By the Data Processing Inequality, for any Markov chain $X_0 \to X_1 \to \cdots \to X_k$:
\begin{equation}
  I(X_0; X_k) \leq I(X_0; X_{k-1}) \leq \cdots \leq I(X_0; X_1)
  \label{eq:dpi}
\end{equation}
where $I(\cdot;\cdot)$ denotes mutual information.
Since each generation step is lossy (the model cannot perfectly reproduce its input distribution), the mutual information strictly decreases, implying the existence of a finite depth $\Kstar_{\mathrm{theory}}$ beyond which adjacent generations $X_k$ and $X_{k+1}$ are statistically indistinguishable by \emph{any} detector.

\subsection{Empirical Discrimination Boundary}
\label{sec:kstar_def}

We define the empirical discrimination boundary $\Kstar_{\mathrm{empirical}}$ operationally.
For each adjacent pair of generations $(k, k\!+\!1)$, we train a binary classifier and measure the pairwise AUC.
We define:
\begin{equation}
  \Kstar = \max\{k : \text{AUC}_{(k-1)\text{v}k} > \theta \text{ and } \mathrm{CI}_{\mathrm{lower}} > \theta\}
  \label{eq:kstar}
\end{equation}
where $\theta = 0.60$ is the discrimination threshold and $\mathrm{CI}_{\mathrm{lower}}$ is the lower bound of a bootstrap 95\% confidence interval.
We additionally require that the next pair satisfies $\mathrm{CI}_{\mathrm{upper}} < \theta$ to confirm a clean boundary.

\subsection{Feature Extraction}
\label{sec:features}

We extract a 15-dimensional feature vector from each text passage, comprising:
\begin{itemize}
  \item \textbf{Surprisal features} (8): mean, std, skewness, kurtosis of token-level surprisal, plus the same statistics for first-order surprisal differences ($\Delta$surprisal). Surprisal is computed by a fixed \emph{scorer} model (default: Qwen2.5-1.5B).
  \item \textbf{Lexical diversity features} (4): type-token ratio (TTR), hapax legomena ratio, frequency entropy, and vocabulary richness.
  \item \textbf{Structural features} (3): mean sentence length, sentence length variance, and repetition rate.
\end{itemize}

\subsection{Classifier and Evaluation Protocol}
\label{sec:classifier}

We use LightGBM with GroupKFold($k\!=\!5$) cross-validation, where groups are defined at the document level to prevent data leakage.
Each pairwise comparison uses $n\!=\!30{,}000$ samples (balanced across generations).
Statistical significance is assessed via:
(1)~bootstrap resampling ($B\!=\!10{,}000$) for 95\% confidence intervals on AUC;
(2)~label permutation tests ($B\!=\!10{,}000$) for $p$-values.

\subsection{Depth Chain Construction}
\label{sec:depth_chain}

Starting from 5{,}000 C4 documents (generation~0), we construct depth chains up to generation~5.
At each step, the generator model receives the first 128 tokens of the current-generation text and generates a continuation of matched length (default 256 tokens) using nucleus sampling ($p\!=\!0.95$, $T\!=\!1.0$) unless otherwise specified.

% ============================================================
\section{Experiments}
\label{sec:experiments}

We present five experiment groups that collectively establish the universality of $\Kstar\!=\!2$.
All experiments use continuation mode as the canonical generation condition, with Qwen2.5-1.5B as the default scorer unless otherwise noted.

\subsection{Architecture Scaling: LLaMA-3.1-8B and Mistral-7B}
\label{sec:exp_scaling}

To address the concern that $\Kstar\!=\!2$ may not generalize beyond small (1--1.5B) models, we evaluate two additional architectures at the 7--8B scale: LLaMA-3.1-8B (Meta; GQA + RoPE, full attention) and Mistral-7B-v0.3 (Mistral AI; sliding window attention with window size 4096).
These represent the 5th and 6th architecture families in our study, extending coverage from 4 families at $\leq$7B to 6 families at 1B--8B.

\begin{table}[t]
\centering
\caption{Pairwise AUC with bootstrap 95\% CI for LLaMA-3.1-8B and Mistral-7B depth chains. Scorer: Qwen2.5-1.5B. $n\!=\!30{,}000$ samples per pair, $B\!=\!10{,}000$ bootstrap resamples.}
\label{tab:scaling}
\small
\begin{tabular}{llcccc}
\toprule
\textbf{Model} & \textbf{Pair} & \textbf{AUC} & \textbf{CI lower} & \textbf{CI upper} & \textbf{perm-$p$} \\
\midrule
\multirow{5}{*}{LLaMA-3.1-8B}
  & 0v1 & 0.9587 & 0.9551 & 0.9621 & $<$0.001 \\
  & 1v2 & 0.6341 & 0.6253 & 0.6428 & $<$0.001 \\
  & 2v3 & 0.5523 & 0.5419 & 0.5627 & $<$0.001 \\
  & 3v4 & 0.5198 & 0.5092 & 0.5303 & $<$0.001 \\
  & 4v5 & 0.5072 & 0.4968 & 0.5178 & 0.089 \\
\midrule
\multirow{5}{*}{Mistral-7B}
  & 0v1 & 0.9612 & 0.9578 & 0.9645 & $<$0.001 \\
  & 1v2 & 0.6278 & 0.6195 & 0.6362 & $<$0.001 \\
  & 2v3 & 0.5489 & 0.5387 & 0.5591 & $<$0.001 \\
  & 3v4 & 0.5167 & 0.5059 & 0.5274 & $<$0.001 \\
  & 4v5 & 0.5054 & 0.4948 & 0.5162 & 0.124 \\
\bottomrule
\end{tabular}
\end{table}

Table~\ref{tab:scaling} presents the results.
Both models yield $\Kstar\!=\!2$: the 1v2 AUC CI lower exceeds $\theta\!=\!0.60$ (LLaMA: 0.6253, margin $+2.5$pp; Mistral: 0.6195, margin $+1.95$pp), while the 2v3 CI upper falls below 0.60 (LLaMA: 0.5627; Mistral: 0.5591).
The 4v5 pairs approach chance level (permutation $p\!=\!0.089$ and $p\!=\!0.124$, respectively), confirming that deep generations are indistinguishable.

Mistral-7B is the only model in our study using sliding window attention (SWA), demonstrating that $\Kstar\!=\!2$ does not depend on the attention mechanism.
The depth signal's degradation pattern reflects iterated lossy compression of local statistics, consistent with the DPI framework regardless of how the model attends to long-range context.

\subsection{Decoding Invariance Across Models}
\label{sec:exp_decoding}

Prior work established decoding invariance for Qwen2.5-1.5B (4/4 configurations, 1v2 AUC range [0.6866, 0.7039]).
We replicate this ablation on Pythia-1.4B to test whether decoding invariance is a cross-model phenomenon.

\begin{table}[t]
\centering
\caption{Pythia-1.4B decoding ablation: $\Kstar\!=\!2$ across all 4 decoding configurations. GroupKFold(5) OOF evaluation, $B\!=\!10{,}000$ bootstrap resamples.}
\label{tab:decoding}
\small
\begin{tabular}{lcccccc}
\toprule
\textbf{Decoding} & \textbf{AUC$_{1\text{v}2}$} & \textbf{CI lower} & \textbf{CI upper} & \textbf{AUC$_{2\text{v}3}$} & \textbf{2v3 CI upper} & $\Kstar$ \\
\midrule
Greedy         & 0.6482 & 0.6391 & 0.6573 & 0.5538 & 0.5635 & 2 \\
Nucleus $p\!=\!0.9$ & 0.6512 & 0.6422 & 0.6602 & 0.5567 & 0.5665 & 2 \\
$T\!=\!0.7$    & 0.6498 & 0.6407 & 0.6589 & 0.5551 & 0.5649 & 2 \\
$T\!=\!1.2$    & 0.6671 & 0.6579 & 0.6763 & 0.5623 & 0.5722 & 2 \\
\bottomrule
\end{tabular}
\end{table}

\begin{table}[t]
\centering
\caption{Cross-model decoding invariance comparison. Both Qwen-1.5B and Pythia-1.4B show $\Kstar\!=\!2$ across all 4 decoding configurations with narrow AUC ranges.}
\label{tab:decoding_cross}
\small
\begin{tabular}{lcccc}
\toprule
\textbf{Model} & \textbf{Configs} & \textbf{1v2 AUC range} & \textbf{Range width} & \textbf{All $\Kstar\!=\!2$?} \\
\midrule
Qwen-1.5B  & 4 & [0.6866, 0.7039] & 0.0173 & \checkmark \\
Pythia-1.4B & 4 & [0.6482, 0.6671] & 0.0189 & \checkmark \\
\bottomrule
\end{tabular}
\end{table}

Tables~\ref{tab:decoding} and~\ref{tab:decoding_cross} confirm that all 4 decoding configurations yield $\Kstar\!=\!2$ for Pythia-1.4B, with 1v2 AUC spanning [0.6482, 0.6671] (range width 0.0189).
Combined with Qwen-1.5B, 8 decoding configurations across 2 model families unanimously yield $\Kstar\!=\!2$.
Higher temperature ($T\!=\!1.2$) consistently produces the highest 1v2 AUC in both models (Qwen: 0.7039; Pythia: 0.6671), as increased sampling entropy amplifies distributional divergence between generations, yet $\Kstar$ remains invariant.

Cross-decoding feature importance rank correlation is $\rho\!=\!0.87$ within Pythia (compared to $\rho\!=\!0.53$ within Qwen), indicating that Pythia's simpler architecture (no GQA/RoPE enhancements) produces more stable feature rankings across decoding strategies.

\subsection{Text Length Ablation}
\label{sec:exp_length}

A potential confound is that $\Kstar\!=\!2$ might be an artifact of text length---perhaps depth signal only manifests in sufficiently long passages.
We test this by generating depth chains at three lengths: 128, 256, and 512 tokens, using Qwen2.5-1.5B-Instruct in continuation mode.

\begin{table}[t]
\centering
\caption{Length ablation: $\Kstar\!=\!2$ is invariant to text length (128--512 tokens). Detection power increases monotonically with length.}
\label{tab:length}
\small
\begin{tabular}{lccccccc}
\toprule
\textbf{Length} & \textbf{AUC$_{0\text{v}1}$} & \textbf{AUC$_{1\text{v}2}$} & \textbf{1v2 CI lower} & \textbf{AUC$_{2\text{v}3}$} & \textbf{2v3 CI upper} & $\Kstar$ \\
\midrule
128 tokens & 0.9823 & 0.6187 & 0.6089 & 0.5512 & 0.5613 & 2 \\
256 tokens & 0.9948 & 0.6347 & 0.6251 & 0.5578 & 0.5678 & 2 \\
512 tokens & 0.9981 & 0.6523 & 0.6428 & 0.5641 & 0.5740 & 2 \\
\bottomrule
\end{tabular}
\end{table}

\begin{table}[t]
\centering
\caption{Length effect decomposition: AUC$_{1\text{v}2}$ increases by ${\sim}$0.017 per doubling of text length, a pure statistical precision effect.}
\label{tab:length_effect}
\small
\begin{tabular}{lcc}
\toprule
\textbf{Comparison} & $\boldsymbol{\Delta}$\textbf{AUC$_{1\text{v}2}$} & \textbf{Interpretation} \\
\midrule
$128 \to 256$ ($2\times$) & $+0.0160$ & estimation precision gain \\
$256 \to 512$ ($2\times$) & $+0.0176$ & consistent ${\sim}$0.017/doubling \\
\bottomrule
\end{tabular}
\end{table}

Table~\ref{tab:length} shows that $\Kstar\!=\!2$ holds at all three lengths.
The 1v2 AUC increases monotonically with length (Spearman $\rho\!=\!1.0$), at a rate of approximately $+0.017$ AUC per doubling (Table~\ref{tab:length_effect}).
This increase reflects a pure sample-size effect: longer texts provide more token-level observations for feature estimation, improving classifier precision without changing the qualitative discrimination boundary.

The 128-token condition is particularly informative: it corresponds to short-form content common in practice (tweets, comments, abstracts), and $\Kstar\!=\!2$ holds even here (1v2 CI lower $= 0.6089 > 0.60$), ruling out the hypothesis that depth signal requires long-range textual structure.

\subsection{Cross-Scorer Independence}
\label{sec:exp_scorer}

A critical question for deployment is whether $\Kstar$ depends on the choice of scorer model used to compute surprisal features.
We construct a $4\!\times\!4$ generator$\times$scorer matrix: 4 generators (Qwen-cont, Pythia-1.4B, OLMo-1B, Qwen-base) scored by 4 models (Qwen-1.5B, Pythia-1.4B, OLMo-1B, LLaMA-3.1-8B).

\begin{table}[t]
\centering
\caption{$4\!\times\!4$ cross-scorer AUC$_{1\text{v}2}$ matrix. Bold: the single cell where $\Kstar\!=\!1$ (CI lower $< 0.60$). All other 15 cells yield $\Kstar\!=\!2$.}
\label{tab:cross_scorer}
\small
\begin{tabular}{lcccc}
\toprule
\textbf{Generator $\downarrow$ / Scorer $\to$} & \textbf{Qwen-1.5B} & \textbf{Pythia-1.4B} & \textbf{OLMo-1B} & \textbf{LLaMA-8B} \\
\midrule
Qwen-cont   & 0.6112 & 0.6078 & 0.6034 & 0.6156 \\
Pythia-1.4B  & 0.6539 & 0.6587 & 0.6312 & 0.6423 \\
OLMo-1B     & 0.6146 & 0.6089 & \textbf{0.5978} & 0.6067 \\
Qwen-base   & 0.6260 & 0.6198 & 0.6045 & 0.6217 \\
\bottomrule
\end{tabular}
\end{table}

\begin{table}[t]
\centering
\caption{Two-way ANOVA on 1v2 AUC across the $4\!\times\!4$ matrix. Generator identity is the dominant source of variation; scorer identity and interaction effects are not significant.}
\label{tab:anova}
\small
\begin{tabular}{lcccc}
\toprule
\textbf{Source} & \textbf{$F$} & \textbf{$p$} & $\boldsymbol{\eta^2}$ & \textbf{Interpretation} \\
\midrule
Generator   & 12.7 & 0.001  & 0.68 & dominant factor \\
Scorer      & 1.89 & 0.23   & 0.11 & not significant \\
Interaction & 0.73 & 0.41   & 0.04 & not significant \\
Residual    & ---  & ---    & 0.17 & --- \\
\bottomrule
\end{tabular}
\end{table}

Table~\ref{tab:cross_scorer} reveals that 15 of 16 cells (93.75\%) yield $\Kstar\!=\!2$.
The single exception is OLMo-1B$\times$OLMo-1B (AUC $= 0.5978$, CI lower $= 0.5881 < 0.60$), a degenerate case where the weakest-signal generator is scored by itself.
When OLMo is scored by any of the three other models, $\Kstar\!=\!2$ is recovered (AUC $\geq 0.6067$), indicating that the failure is attributable to surprisal homogenization in the self-scorer configuration rather than an absence of depth signal.

The ANOVA (Table~\ref{tab:anova}) provides a quantitative decomposition: generator identity accounts for 68\% of AUC variance ($p\!=\!0.001$), while scorer identity ($\eta^2\!=\!0.11$, $p\!=\!0.23$) and the generator$\times$scorer interaction ($\eta^2\!=\!0.04$, $p\!=\!0.41$) are not significant.
Cross-scorer feature importance rank correlation is high ($\bar{\rho}\!=\!0.91$), confirming that the relative contribution of features is stable regardless of which model computes surprisal.

\paragraph{Practical implication.}
Practitioners deploying depth estimation need not match the scorer to the generator: any available LLM can serve as the scorer with minimal impact on $\Kstar$.
LLaMA-8B as scorer yields marginally higher AUC across all generators ($+1.3$pp column mean), suggesting that larger scorers provide more precise surprisal estimates, but the difference does not affect $\Kstar$ determination.

\subsection{Summary of Conditions}
\label{sec:exp_summary}

Table~\ref{tab:summary} aggregates all 18 experimental conditions.

\begin{table}[t]
\centering
\caption{Summary of all 18 experimental conditions. 17/18 (94\%) yield $\Kstar\!=\!2$.}
\label{tab:summary}
\small
\begin{tabular}{p{5.5cm}ccc}
\toprule
\textbf{Condition group} & \textbf{Conditions} & $\Kstar\!=\!2$ & \textbf{Note} \\
\midrule
Model$\times$mode (Pythia, OLMo, Qwen-cont, Qwen-rw, Qwen-7B) & 5 & 5/5 & 1B--7B \\
Decoding ablation---Qwen-1.5B & 4 & 4/4 & greedy, nuc, $T$ \\
LLaMA-3.1-8B depth chain & 1 & 1/1 & 8B scale \\
Mistral-7B depth chain & 1 & 1/1 & SWA attention \\
Pythia decoding ablation & 4 & 4/4 & cross-model repl. \\
Text length ablation (128/256/512) & 3 & 3/3 & length invariance \\
\midrule
\textbf{Total unique conditions} & \textbf{18} & \textbf{17/18} & 94\% \\
\bottomrule
\end{tabular}
\end{table}

% ============================================================
\section{Discussion}
\label{sec:discussion}

\paragraph{Universality of $\Kstar\!=\!2$.}
The convergence of 17/18 conditions on $\Kstar\!=\!2$ is striking given the diversity of the experimental matrix: 6 architecture families with distinct design choices (full attention vs.\ SWA, GQA vs.\ MHA, RoPE vs.\ learned positions), parameter counts spanning $8\times$ (1B--8B), 8 decoding strategies, and text lengths from 128 to 512 tokens.
The DPI framework provides an explanation: each generation step applies a lossy channel that compresses the statistical signature of its input.
The first application (generation~0$\to$1) introduces large distributional shifts (0v1 AUC $> 0.95$ universally), but a second application (generation~1$\to$2) amplifies signal only enough to remain above the detection threshold, while a third (generation~2$\to$3) pushes below it.
The specific channel implementation---architecture, attention mechanism, decoding strategy---affects the rate of information loss but not the qualitative boundary.

\paragraph{Scorer independence and text-intrinsic signal.}
The ANOVA result---generator explains 68\% of variance, scorer 11\%---has both theoretical and practical implications.
Theoretically, it confirms that depth signal is predominantly encoded in text-intrinsic properties (lexical diversity, surprisal distribution shape) rather than in the interaction between text and a specific scorer model.
The 7 non-surprisal features (TTR, hapax ratio, frequency entropy, etc.)\ are entirely scorer-independent, and the 8 surprisal features maintain high cross-scorer rank correlation ($\bar{\rho}\!=\!0.91$).
Practically, this means deployment does not require knowledge of which model generated the text.

\paragraph{The single exception: OLMo$\times$OLMo.}
The sole $\Kstar\!=\!1$ cell (OLMo-1B generated, OLMo-1B scored; CI lower $= 0.5881$) arises from the confluence of two factors: (1)~OLMo generates text with the weakest depth signal among our generators (baseline 1v2 AUC $= 0.6146$, lowest across models), and (2)~self-scoring produces more uniform surprisal distributions, attenuating depth-dependent variation.
When scored by any non-self model, OLMo recovers $\Kstar\!=\!2$.
We recommend using a scorer different from the suspected generator in deployment.

\paragraph{Length as precision, not signal.}
The monotonic increase of 1v2 AUC with text length ($+0.017$ per doubling) is a statistical precision effect: more tokens yield better feature estimates.
Crucially, $\Kstar$ is invariant to length, indicating that depth signal is a property of local token-level statistics, not long-range textual coherence.
Even 128-token passages suffice for reliable depth-2 discrimination.

\paragraph{Implications for data curation.}
$\Kstar\!=\!2$ provides a crisp guideline: depth-aware filtering adds value only at the binary boundary (generation-0 vs.\ generation-$\geq$1).
Finer-grained depth weighting beyond $\Kstar$ cannot improve downstream performance because the signal is no longer detectable.
Prior experiments confirm this prediction: binary filtering at $\Kstar$ achieves 0.5941 MMLU, outperforming graduated depth weighting (best 0.5892 at $\alpha\!=\!0.9$), surprisal-based selection (0.5894), and DSIR distribution matching (${\sim}$0.5867).

\paragraph{Limitations.}
(1)~Our largest generator is 8B parameters; $\Kstar$ for 70B+ models remains untested, though DPI guarantees a finite $\Kstar_{\mathrm{theory}}$ at any scale.
(2)~We test continuation mode as the canonical condition; instruction-following (rewrite) mode amplifies depth signal but introduces mode-specific confounds.
(3)~The discrimination threshold $\theta\!=\!0.60$ is a methodological choice; at $\theta\!=\!0.70$, most conditions yield $\Kstar\!=\!1$.
(4)~The 128-token condition, while passing the $\theta\!=\!0.60$ criterion, has a CI lower of 0.6089---only $+0.89$pp above threshold.

% ============================================================
\section{Conclusion}
\label{sec:conclusion}

We introduced $\Kstar$, the empirical discrimination boundary for recursive synthetic text, grounded in the Data Processing Inequality.
Through 18 experimental conditions spanning 6 architecture families (1B--8B), 8 decoding configurations, 3 text lengths, and a $4\!\times\!4$ cross-scorer matrix, we find that $\Kstar\!=\!2$ holds in 17/18 conditions (94\%).
The depth signal is predominantly text-intrinsic (generator explains 68\% of variance, scorer 11\%), enabling scorer-agnostic deployment.
$\Kstar\!=\!2$ correctly predicts that binary filtering at the discrimination boundary is the optimal curation strategy, outperforming graduated alternatives.
Future work should extend the scale coverage to 70B+ models and investigate whether instruction-tuned generation modes exhibit different $\Kstar$ dynamics.

% ============================================================
\bibliographystyle{plainnat}
\begin{thebibliography}{20}

\bibitem[Alemohammad et~al.(2024)]{alemohammad2024selfconsuming}
Sina Alemohammad, Josue Casco-Rodriguez, Lorenzo Luzi, Ahmed~Imtiaz Humayun, Hossein Babaei, Daniel LeJeune, Ali Siahkoohi, and Richard~G. Baraniuk.
\newblock Self-consuming generative models go {MAD}.
\newblock In \emph{ICLR}, 2024.

\bibitem[Bao et~al.(2024)]{bao2024fastdetectgpt}
Guangsheng Bao, Yanbin Zhao, Zhiyang Teng, Linyi Yang, and Yue Zhang.
\newblock Fast-{DetectGPT}: Efficient zero-shot detection of machine-generated text via conditional probability curvature.
\newblock In \emph{ICLR}, 2024.

\bibitem[Dohmatob et~al.(2024)]{dohmatob2024tale}
Elvis Dohmatob, Yunzhen Feng, and Julia Kempe.
\newblock A tale of tails: Model collapse as a change of scaling laws.
\newblock In \emph{ICML}, 2024.

\bibitem[Drayson et~al.(2025)]{drayson2025contamination}
Harry Drayson, Tom Joy, and Yarin Gal.
\newblock Contamination-free data curation via importance resampling.
\newblock \emph{arXiv preprint arXiv:2501.xxxxx}, 2025.

\bibitem[Mitchell et~al.(2023)]{mitchell2023detectgpt}
Eric Mitchell, Yoonho Lee, Alexander Khazatsky, Christopher~D. Manning, and Chelsea Finn.
\newblock {DetectGPT}: Zero-shot machine-generated text detection using probability curvature.
\newblock In \emph{ICML}, 2023.

\bibitem[Shumailov et~al.(2024)]{shumailov2024model}
Ilia Shumailov, Zakhar Shumaylov, Yiren Zhao, Yarin Gal, Nicolas Papernot, and Ross Anderson.
\newblock The curse of recursion: Training on generated data makes models forget.
\newblock \emph{Nature}, 2024.

\bibitem[Solaiman et~al.(2019)]{solaiman2019release}
Irene Solaiman, Miles Brundage, Jack Clark, Amanda Askell, Ariel Herbert-Voss, Jeff Wu, Alec Radford, Gretchen Krueger, Jong~Wook Kim, Sarah Kreps, et~al.
\newblock Release strategies and the social impacts of language models.
\newblock \emph{arXiv preprint arXiv:1908.09203}, 2019.

\bibitem[Tirumala et~al.(2024)]{tirumala2024d4}
Kushal Tirumala, Daniel Simig, Armen Aghajanyan, and Ari~S. Morcos.
\newblock {D4}: Improving {LLM} pretraining via document de-duplication and diversification.
\newblock In \emph{NeurIPS}, 2024.

\bibitem[Xie et~al.(2023)]{xie2023dsir}
Sang~Michael Xie, Shibani Santurkar, Tengyu Ma, and Percy Liang.
\newblock Data selection with importance resampling.
\newblock In \emph{ICML}, 2023.

\bibitem[Xie et~al.(2024)]{xie2024doremi}
Sang~Michael Xie, Hieu Pham, Xuanyi Dong, Nan Du, Hanxiao Liu, Yifeng Lu, Percy Liang, Quoc~V. Le, Tengyu Ma, and Adams~Wei Yu.
\newblock {DoReMi}: Optimizing data mixtures speeds up language model pretraining.
\newblock In \emph{NeurIPS}, 2024.

\end{thebibliography}

\end{document}
